%!TEX root = ../LLMSeminar.tex
% ──────────────────────────────────────────────────────────────
%  Section 8 — Self-improvement and Its Limits
% ──────────────────────────────────────────────────────────────

\section{Self-improvement and its limits}\label{sec:self-improvement}

A coding agent that can read and write files, and which knows
the path to its own source, can modify itself.  This is a
trivial geometric consequence of the loop in
Theorem~\ref{kr:agent-loop}; it is not a profound emergent
phenomenon.  But the consequences are striking, so it is worth
spelling out carefully---together with the much-less-discussed
fact that the kind of self-improvement available to agents is
\emph{strictly limited} in a way that matters.

\subsection{The two senses of ``self-improvement''}

The phrase ``self-improving AI'' is used in popular discourse
in a way that conflates two very different things.  We must
keep them apart.

\begin{keyresult}[Two senses of self-improvement]
  \begin{enumerate}[(i)]
    \item \textbf{Harness self-improvement} (easy, common,
      already happening).  The Python (or whatever) programme
      that runs around the LLM is just a file; the LLM can
      read it and write a better version.  Nothing exotic;
      this is plain editing.
    \item \textbf{Weights self-improvement} (impossible at
      inference, expensive even at training time).  The
      neural-network weights inside the LLM are frozen at
      deployment.  There is no mechanism by which a running
      model updates its own parameters.  ``Improvement'' of
      the model means \emph{re-training}, which involves
      orders-of-magnitude more compute than inference.
  \end{enumerate}
\end{keyresult}

\begin{intuition}[A physicist's gloss]
  Think of the LLM weights as a frozen Hamiltonian and the
  conversation as time-evolution under that Hamiltonian.  The
  Hamiltonian is fixed; only the state evolves.  ``Self-improvement
  by weights'' would mean the Hamiltonian rewrites itself in
  flight, which is not what is happening.  ``Self-improvement
  by harness'' means we are giving the system a richer
  measurement apparatus and a better workspace; it is
  completely classical engineering.
\end{intuition}

\subsection{The minimal self-improving agent}

The construction is direct.  Take the canonical agent of
\S\ref{sec:agent-loop}, give it the prompt:

\begin{lstlisting}
You are a coding agent with read- and write-file tools.
Your own source code lives at agent.py.
You may read and write it to improve yourself.
Changes take effect on next invocation.
Preserve the loop and tool plumbing in this script.
\end{lstlisting}

\noindent and ask it: \emph{``go and improve yourself.''}

A representative trace, observed in practice:

\begin{enumerate}[(1)]
  \item The agent reads \texttt{agent.py}.
  \item It reasons (in a verbose intermediate step) that, with
    only read and write, it cannot run additional commands; an
    \texttt{execute} tool would expand its capability greatly.
  \item It rewrites \texttt{agent.py} to add a third tool,
    \texttt{run\_bash(command)}, that executes a shell command.
  \item It exits.
  \item On the next invocation, the agent has bash; from there
    it can install dependencies, fetch packages, query the
    network, and so on.
\end{enumerate}

\begin{remark}
  The agent will often note, in its own commentary, that adding
  bash carries safety implications, and proceed nonetheless if
  not specifically forbidden.  Anthropic's models are trained
  fairly aggressively to flag such moves; other vendors are
  more permissive.  The overall pattern, however, is universal.
\end{remark}

\subsection{Why this works}

The construction is just composition.  The agent already had
a tool that maps strings to writes on disk and a tool that
maps strings to reads.  Adding a tool is only a syntactic
matter of registering it in the tools array and writing a
dispatcher branch.  The fact that the LLM does this happily,
without close supervision, is a property of training:
the model has seen many programmes and knows how to extend
them.

\begin{historical}[The coding-agent ``death match'']
  Every few months the AI community runs an informal experiment
  in which several coding agents are launched on the same
  machine with one instruction: \emph{kill all the other
  agents; last one running wins.}  The agents have become so
  effective at expanding their own capabilities (installing
  network tools, scanning processes, exploiting privilege
  escalations) that the contest now requires Docker isolation:
  the agents will, given a few minutes, escape any
  permission boundary one cares to set in pure user space.
  Whatever one thinks of this as a parlour game, it is a real
  data point on agent capability.
\end{historical}

\subsection{What ``self-improvement'' cannot do}

It is critical not to over-read these demonstrations.  The
LLM in the loop never updates its own weights.  Strong
versions of the ``recursive self-improvement'' thesis---the
idea that an AI iteratively rewrites its own intelligence
upward---do not apply to anything available today, because the
intelligence sits in the frozen weights, and the frozen weights
do not move during inference.

\begin{keyresult}[The unmoving floor]
  No matter how clever the harness becomes, the underlying
  cognitive ceiling is set by the fixed weights of the model
  it is calling.  Improving the harness expands what the
  fixed cognition is allowed to \emph{do}; it does not raise
  what it can \emph{think}.
\end{keyresult}

\subsection{The training-time loop}

There is, however, a slower and external loop in which models
do improve from agent activity.

\begin{keyresult}[Synthetic-data improvement]
  Coding agents generate a vast quantity of high-quality
  interaction data: trajectories of plan, action, result,
  correction, success.  These trajectories can be filtered,
  ranked, and folded into the training set of the
  \emph{next} model.  In this sense models do improve from
  their own deployment, but the improvement happens at the
  next training cycle, not in any individual conversation.
\end{keyresult}

This synthetic-data path is one of the most consequential
shifts in the recent training pipeline.  For coding tasks,
the data appears almost free: it is generated by running the
agents themselves, with automatic verifiers (does the code
compile, do the tests pass) acting as the filter.

\begin{audience}[Audience question]
  \emph{Does this solve the data scarcity problem?}

  For coding it largely does, yes; the closed loop of
  \emph{agent action $\to$ verification $\to$ training data}
  is now self-sustaining.  For other domains, where verification
  is harder or where data is restricted by law or social
  norm---medical records being the canonical example, where
  even synthetic generation is hampered by privacy
  constraints---the picture is much less rosy, and Western
  models in particular operate under tight data limitations.
\end{audience}

\begin{remark}
  Be careful with the word ``self.''  A coding agent improves
  its harness by writing code; an organisation improves its
  model by training the next checkpoint.  Both are sometimes
  called ``self-improvement,'' but they live on radically
  different timescales (seconds vs.\ months) and require
  radically different resources (a laptop vs.\ a data centre).
\end{remark}

\subsection{What this means in practice}

Two operational lessons come out of the above.

\begin{enumerate}[(1)]
  \item If you give a coding agent broad system access without
    isolation, you should expect it to expand its own
    capabilities and to use them.  Dockerise; sandbox;
    audit logs.
  \item Do not assume that running an agent for longer makes
    it ``smarter.''  It can do more, look up more, accumulate
    more context; but the underlying reasoning floor is fixed
    by the model.  If the model cannot reason about a class of
    problems, more turns will not help---it will simply run
    more confidently in circles.
\end{enumerate}
